\documentclass[a4paper,12pt]{article}

% ------------------------------------------------------------
% Кодировка, язык, математика
% ------------------------------------------------------------
\usepackage[T2A]{fontenc}
\usepackage[utf8]{inputenc}
\usepackage[russian]{babel}
\usepackage{amsmath,amssymb,mathtools}
\usepackage{icomma}

% ------------------------------------------------------------
% Типографика и страница
% ------------------------------------------------------------
\usepackage{PTSans}
\renewcommand{\familydefault}{\sfdefault}
\usepackage{microtype}
\usepackage{setspace}
\onehalfspacing
\usepackage{parskip}
\usepackage[
  a4paper,
  left=20mm,
  right=20mm,
  top=20mm,
  bottom=22mm
]{geometry}

% ------------------------------------------------------------
% Графика
% ------------------------------------------------------------
\usepackage{tikz}
\usetikzlibrary{calc,arrows.meta,positioning,shapes.geometric}

% ------------------------------------------------------------
% Таблицы, списки, служебные пакеты
% ------------------------------------------------------------
\usepackage{tabularx,booktabs}
\usepackage{enumitem}
\usepackage{needspace}
\usepackage{xcolor}
\usepackage{hyperref}
\usepackage{fancyhdr}

% ------------------------------------------------------------
% Палитра
% ------------------------------------------------------------
\definecolor{accent}{HTML}{008080}
\definecolor{darkaccent}{HTML}{004D4D}
\definecolor{textgray}{HTML}{2F3333}
\definecolor{lightgray}{HTML}{D9E2E2}

\color{textgray}

% ------------------------------------------------------------
% PDF-метаданные
% ------------------------------------------------------------
\hypersetup{
  hidelinks,
  pdfauthor={Лёвин Артём Александрович},
  pdftitle={Чек-лист: линейные неравенства, системы и задачи по формулам}
}

% ------------------------------------------------------------
% Колонтитулы
% ------------------------------------------------------------
\pagestyle{fancy}
\fancyhf{}
\fancyhead[L]{\small Ярослав Верещагин}
\fancyhead[R]{\small Подготовка к ОГЭ}
\fancyfoot[C]{\small\thepage}
\setlength{\headheight}{15pt}
\renewcommand{\headrulewidth}{0.3pt}
\renewcommand{\footrulewidth}{0pt}

% ------------------------------------------------------------
% Служебные команды оформления
% ------------------------------------------------------------
\newcommand{\topic}[1]{%
  \Needspace{9\baselineskip}%
  \par\bigskip
  {\Large\bfseries\color{darkaccent}#1}\par
  \vspace{1mm}
  {\color{accent!55}\rule{\linewidth}{0.5pt}}\par
  \medskip
}

\newcommand{\keyidea}[1]{%
  \par\smallskip
  \noindent{\color{darkaccent}\textbf{Ключевая идея.}} #1\par\smallskip
}

\newcommand{\warning}[1]{%
  \par\smallskip
  \noindent{\color{darkaccent}\textbf{Обратите внимание.}} #1\par\smallskip
}

\newcommand{\exampletitle}[1]{%
  \Needspace{8\baselineskip}%
  \par\medskip
  \noindent\textbf{#1}\par\smallskip
}

% ------------------------------------------------------------
% Стили блок-схемы
% ------------------------------------------------------------
\tikzset{
  flowstart/.style={
    ellipse,
    draw=accent!75,
    line width=0.8pt,
    align=center,
    minimum width=4cm,
    minimum height=9mm,
    font=\small
  },
  flowaction/.style={
    rectangle,
    rounded corners=2mm,
    draw=accent!65,
    line width=0.8pt,
    align=center,
    text width=7cm,
    minimum height=11mm,
    inner sep=3mm,
    font=\small
  },
  flowdecision/.style={
    diamond,
    aspect=2.2,
    draw=darkaccent!80,
    line width=0.8pt,
    align=center,
    text width=3.6cm,
    inner sep=1.5mm,
    font=\small
  },
  flowside/.style={
    rectangle,
    rounded corners=2mm,
    draw=accent!55,
    line width=0.75pt,
    align=center,
    text width=3.5cm,
    minimum height=12mm,
    inner sep=2.5mm,
    font=\small
  },
  flowarrow/.style={
    -{Stealth[length=2.5mm,width=1.8mm]},
    line width=0.75pt,
    draw=textgray!85
  }
}

\begin{document}

% ============================================================
% Титульный лист
% ============================================================
\begin{titlepage}
\thispagestyle{empty}
\begin{center}

\vspace*{2.2cm}

{\color{darkaccent}\Large\bfseries ЧЕК-ЛИСТ}

\vspace{0.7cm}

{\Huge\bfseries
Линейные неравенства, системы\\[2mm]
и задачи по формулам
}

\vspace{0.9cm}

{\large Подготовка к ОГЭ по математике}

\vspace{1.8cm}

{\Large\bfseries Ярослав Верещагин}

\vspace{0.7cm}

{\normalsize
Числовые промежутки \(\bullet\) линейные неравенства \(\bullet\) системы\\
подстановка в формулы \(\bullet\) сложные дроби
}

\vfill

{\large 17 сентября 2026 года}

\vspace{1cm}

{\normalsize
Лёвин Артём Александрович\\[1mm]
\textit{эксклюзивно для Ярослава Верещагина}
}

\vspace{1.6cm}

\begin{tikzpicture}
\draw[accent!45,line width=1.4pt]
(-6.5,0)
.. controls (-4.1,1.0) and (-2.3,-0.75) .. (0,0.05)
.. controls (2.2,0.82) and (4.1,-0.62) .. (6.5,0.08);
\end{tikzpicture}

\end{center}
\end{titlepage}

% ============================================================
% Основной чек-лист
% ============================================================
\topic{1. Как читать неравенство и записывать ответ промежутком}

\keyidea{Неравенство описывает множество чисел, для которых условие верно.}

Например,
\[
x\ge5
\]
читается: «\(x\) не меньше 5». Число \(5\) подходит, поэтому граница включена:
\[
x\in[5; +\infty).
\]

Если
\[
x>5,
\]
то число \(5\) служит границей и в множество решений не входит:
\[
x\in(5; +\infty).
\]

Аналогично,
\[
x\le5 \quad\Longleftrightarrow\quad x\in(-\infty;5],
\]
\[
x<5 \quad\Longleftrightarrow\quad x\in(-\infty;5).
\]

\begin{center}
\begin{tikzpicture}[x=1.15cm,y=1cm]
  \draw[-{Stealth[length=2mm]}] (-0.6,0) -- (6.6,0);
  \foreach \x/\lab in {0/2,1/3,2/4,3/5,4/6,5/7}
    \draw (\x,-0.09)--(\x,0.09) node[below=3pt]{\small \lab};
  \fill[accent] (3,0) circle (2.1pt);
  \draw[accent,line width=2pt,-{Stealth[length=2mm]}]
       (3.05,0.28)--(6.15,0.28);
  \node[above=5pt,text=darkaccent] at (4.55,0.28) {\small \(x\ge5\)};
\end{tikzpicture}
\end{center}

\warning{Для \(\le\) и \(\ge\) граница включена: точка закрашенная, скобка квадратная. Для \(<\) и \(>\) граница исключена: точка выколотая, скобка круглая. У \(\pm\infty\) всегда ставится круглая скобка.}

% ------------------------------------------------------------
\topic{2. Линейные неравенства}

Решение линейного неравенства похоже на решение линейного уравнения:
\begin{enumerate}[leftmargin=8mm,itemsep=1mm]
  \item раскрыть скобки, если они есть;
  \item перенести слагаемые с \(x\) в одну часть, числа --- в другую;
  \item привести подобные;
  \item разделить обе части на коэффициент при \(x\);
  \item записать ответ промежутком.
\end{enumerate}

\keyidea{Если обе части неравенства умножают или делят на отрицательное число, знак неравенства меняют на противоположный.}

\exampletitle{Пример 1. Положительный коэффициент при \(x\)}
\[
4x+1\ge2x+10,
\]
\[
4x-2x\ge10-1,
\]
\[
2x\ge9,
\]
\[
x\ge4{,}5.
\]
Следовательно,
\[
\boxed{x\in[4{,}5; +\infty)}.
\]

\exampletitle{Пример 2. Деление на отрицательное число}
\[
-4x+1>2x+10,
\]
\[
-6x>9.
\]
Делим обе части на \(-6\) и разворачиваем знак:
\[
x<-\frac{3}{2}.
\]
Следовательно,
\[
\boxed{x\in\left(-\infty;-\frac32\right)}.
\]

\warning{При переносе слагаемого через знак неравенства меняется знак самого слагаемого. Сам знак неравенства разворачивается только при умножении или делении обеих частей на отрицательное число.}

% ============================================================
% Отдельная страница блок-схемы
% ============================================================
\clearpage
\thispagestyle{empty}

\begin{center}
{\Large\bfseries\color{darkaccent}Алгоритм: решение линейного неравенства}

\vspace{8mm}

\begin{tikzpicture}[node distance=8mm and 18mm]

\node[flowstart] (start) {Начало};

\node[flowaction,below=of start] (simplify)
{Раскройте скобки и приведите подобные};

\node[flowaction,below=of simplify] (collect)
{Соберите слагаемые с \(x\) в одной части, числа --- в другой};

\node[flowdecision,below=10mm of collect] (remain)
{После упрощения осталось \(x\)?};

\node[flowside,right=8mm of remain] (numeric)
{Получилось числовое неравенство.\\
Если оно истинно: \(\mathbb{R}\).\\
Если ложно: \(\emptyset\)};

\node[flowdecision,below=11mm of remain] (negative)
{Коэффициент при \(x\) отрицательный?};

\node[flowside,left=8mm of negative] (keep)
{Разделите на коэффициент.\\
Знак сохраняется};

\node[flowside,right=8mm of negative] (flip)
{Разделите на коэффициент.\\
Знак разворачивается};

\node[flowaction,below=19mm of negative] (interval)
{Запишите множество решений числовым промежутком};

\node[flowstart,below=of interval] (finish) {Конец};

\draw[flowarrow] (start) -- (simplify);
\draw[flowarrow] (simplify) -- (collect);
\draw[flowarrow] (collect) -- (remain);
\draw[flowarrow] (remain.east) -- node[above,font=\small]{нет} (numeric.west);
\draw[flowarrow] (remain.south) -- node[right,font=\small]{да} (negative.north);
\draw[flowarrow] (negative.west) -- node[above,font=\small]{нет} (keep.east);
\draw[flowarrow] (negative.east) -- node[above,font=\small]{да} (flip.west);
\draw[flowarrow] (keep.south) |- (interval.west);
\draw[flowarrow] (flip.south) |- (interval.east);
\draw[flowarrow] (interval) -- (finish);
\draw[flowarrow] (numeric.south) |- (finish.east);

\end{tikzpicture}
\end{center}

\clearpage

% ------------------------------------------------------------
\topic{3. Системы неравенств}

В системе должны выполняться \textbf{все условия одновременно}. Поэтому после решения каждого неравенства ищут \textbf{пересечение} полученных промежутков.

\exampletitle{Простой пример}
\[
\begin{cases}
x\ge5,\\
x<10.
\end{cases}
\]

На числовой прямой подходят числа от \(5\) включительно до \(10\) исключительно:
\[
\boxed{x\in[5;10)}.
\]

\begin{center}
\begin{tikzpicture}[x=1cm,y=1cm]
  \draw[-{Stealth[length=2mm]}] (-0.8,0) -- (8.3,0);
  \draw (2,-0.12)--(2,0.12) node[below=3pt]{\small 5};
  \draw (7,-0.12)--(7,0.12) node[below=3pt]{\small 10};
  \fill[accent] (2,0) circle (2.2pt);
  \draw[accent,line width=1.7pt,-{Stealth[length=2mm]}]
       (2.05,0.32)--(7.95,0.32);
  \draw[fill=white,draw=darkaccent,line width=1pt]
       (7,0) circle (2.3pt);
  \draw[darkaccent,line width=1.7pt,-{Stealth[length=2mm]}]
       (6.95,-0.32)--(-0.35,-0.32);
  \draw[textgray!75,line width=3pt] (2.05,-0.72)--(6.95,-0.72);
  \node[below=3pt] at (4.5,-0.72) {\small пересечение};
\end{tikzpicture}
\end{center}

\exampletitle{Пример с преобразованиями}
\[
\begin{cases}
2x-3\ge x+2,\\
4x+1<3x+9.
\end{cases}
\]
Первое неравенство даёт
\[
x\ge5,
\]
второе ---
\[
x<8.
\]
Значит,
\[
\boxed{x\in[5;8)}.
\]

\warning{Если промежутки не пересекаются, система не имеет решений. Например, условия \(x\ge6\) и \(x<2{,}3\) одновременно выполнить нельзя, поэтому ответ --- \(\emptyset\). Если одна и та же граница хотя бы в одном обязательном условии исключена, в итоговое пересечение она тоже не входит.}

% ------------------------------------------------------------
\topic{4. Задачи по готовой формуле}

\keyidea{В таких задачах важно правильно сопоставить величины с буквами в формуле, подставить данные и аккуратно выразить неизвестное.}

Полезный порядок действий:
\begin{enumerate}[leftmargin=8mm,itemsep=1mm]
  \item определить, что дано и что требуется найти;
  \item проверить единицы измерения;
  \item подставить данные в формулу;
  \item упростить выражение, по возможности сначала выполнить сокращения;
  \item решить получившееся уравнение;
  \item проверить ограничения и смысл ответа.
\end{enumerate}

\exampletitle{Пример 1. Прямая подстановка}
Пусть
\[
C=6000+400n,
\]
где \(n\) --- количество дополнительных элементов. При \(n=5\):
\[
C=6000+400\cdot5=8000.
\]
Следовательно,
\[
\boxed{C=8000}.
\]

\exampletitle{Пример 2. Выражение неизвестной}
Центростремительное ускорение задаётся формулой
\[
a=\omega^2r.
\]
Если \(a=64\), \(\omega=4\), то
\[
64=4^2r=16r,
\]
поэтому
\[
\boxed{r=4}.
\]

\exampletitle{Пример 3. Формула с дробью}
Площадь ромба вычисляется по формуле
\[
S=\frac12 d_1d_2\sin\alpha.
\]
Пусть \(S=21\), \(d_1=7\), \(d_2=15\). Тогда
\[
21=\frac12\cdot7\cdot15\cdot\sin\alpha.
\]
Умножаем обе части на \(2\):
\[
42=7\cdot15\cdot\sin\alpha.
\]
Отсюда
\[
\sin\alpha=\frac{42}{7\cdot15}
=\frac{6}{15}
=\frac25
=0{,}4.
\]
Следовательно,
\[
\boxed{\sin\alpha=0{,}4}.
\]

\warning{Не спешите перемножать большие числа. Сначала ищите сокращения: это уменьшает объём вычислений и снижает риск арифметической ошибки.}

% ------------------------------------------------------------
\topic{5. Произведение, сумма и внешний множитель}

При умножении произведения на число внешний множитель записывают один раз:
\[
7\left(x\cdot\frac27\right)
=7\cdot x\cdot\frac27
=2x.
\]

При умножении суммы число распределяют по каждому слагаемому:
\[
7(a+b)=7a+7b.
\]

\keyidea{Распределительный закон нужен для суммы или разности. Для уже записанного произведения дополнительное «распределение по множителям» не выполняют.}

% ------------------------------------------------------------
\topic{6. Многоэтажные дроби}

Деление на дробь заменяют умножением на обратную:
\[
\frac{\frac{a}{b}}{\frac{c}{d}}
=
\frac{a}{b}\cdot\frac{d}{c}
=
\frac{ad}{bc},
\]
где
\[
b\ne0,\qquad c\ne0,\qquad d\ne0.
\]

\exampletitle{Пример 1. Трёхэтажная запись}
\[
\frac{3}{\frac57}
=
3\cdot\frac75
=
\boxed{\frac{21}{5}}.
\]

\exampletitle{Пример 2. Четырёхэтажная дробь}
\[
\frac{\frac{7}{12}}{\frac{14}{15}}
=
\frac{7}{12}\cdot\frac{15}{14}.
\]
Сначала сокращаем:
\[
\frac{7}{14}=\frac12,
\qquad
\frac{15}{12}=\frac54.
\]
Поэтому
\[
\frac{7}{12}\cdot\frac{15}{14}
=
\frac12\cdot\frac54
=
\boxed{\frac58}.
\]

\warning{Если обычное число записано как часть сложной дроби, его можно представить дробью со знаменателем \(1\): например, \(3=\frac31\).}

% ------------------------------------------------------------
\topic{7. Квадрат неизвестной в формуле}

Кинетическая энергия задаётся формулой
\[
E=\frac{mv^2}{2}.
\]

Пусть \(E=98\) кДж и \(m=1000\) кг. Сначала переводим энергию в джоули:
\[
98\text{ кДж}=98000\text{ Дж}.
\]
Подставляем:
\[
98000=\frac{1000v^2}{2}.
\]
Умножаем обе части на \(2\) и делим на \(1000\):
\[
v^2=\frac{2\cdot98000}{1000}=196.
\]
Алгебраическое уравнение \(v^2=196\) имеет два корня:
\[
v=\pm14.
\]
Поскольку \(v\) обозначает модуль скорости, берём неотрицательное значение:
\[
\boxed{v=14\text{ м/с}}.
\]

\warning{При извлечении квадратного корня в алгебраическом уравнении нужно учитывать оба знака. В прикладной задаче затем выбирают значения, допустимые по смыслу величины. Делить на массу можно, поскольку физическая масса положительна: \(m>0\).}

% ------------------------------------------------------------
\topic{8. Быстрая самопроверка}

Перед записью ответа проверьте:
\begin{enumerate}[leftmargin=8mm,itemsep=0.6mm,topsep=1mm]
  \item правильно ли переписан знак \(<,>,\le,\ge\);
  \item развёрнут ли знак после деления на отрицательное число;
  \item соответствует ли тип скобки строгой или нестрогой границе;
  \item в системе найдено именно пересечение промежутков;
  \item при отсутствии пересечения записано \(\emptyset\);
  \item данные подставлены в правильные места формулы;
  \item единицы измерения согласованы;
  \item при делении дробей вторая дробь перевёрнута;
  \item выполнены возможные сокращения до больших вычислений;
  \item учтены ограничения и смысл полученного ответа.
\end{enumerate}

% ============================================================
% Тренировочный блок
% ============================================================
\clearpage

\begin{center}
{\Large\bfseries\color{darkaccent}Тренировочный блок}

\vspace{2mm}
{\small Путь А \(\bullet\) 5 заданий \(\bullet\) от простого к более содержательному}
\end{center}

\vspace{6mm}

\textbf{1.} Запишите множество решений неравенства числовым промежутком:
\[
x\le-2.
\]

\vspace{6mm}

\textbf{2.} Решите неравенство и запишите ответ промежутком:
\[
-5x+8>2x-13.
\]

\vspace{6mm}

\textbf{3.} Решите систему неравенств:
\[
\begin{cases}
2x-3\ge x+2,\\
4x+1<3x+9.
\end{cases}
\]

\vspace{6mm}

\textbf{4.} Вычислите:
\[
\frac{\frac{7}{12}}{\frac{14}{15}}.
\]

\vspace{6mm}

\textbf{5.} Кинетическая энергия тела вычисляется по формуле
\[
E=\frac{mv^2}{2}.
\]
Энергия тела равна \(90\) кДж, масса --- \(800\) кг. Найдите скорость \(v\) в метрах в секунду.

\vfill

\begin{center}
\small
Сначала выполняйте смысловые преобразования и сокращения, затем вычисляйте.
\end{center}

% ============================================================
% Ответы
% ============================================================
\clearpage

\begin{center}
{\Large\bfseries\color{darkaccent}Ответы}
\end{center}

\vspace{6mm}

\renewcommand{\arraystretch}{1.45}
\begin{table}[h!]
\centering
\begin{tabularx}{0.88\linewidth}{
  >{\centering\arraybackslash}p{16mm}
  >{\centering\arraybackslash}X
}
\toprule
\textbf{№} & \textbf{Ответ} \\
\midrule
1 & \(x\in(-\infty;-2]\) \\
2 & \(x\in(-\infty;3)\) \\
3 & \(x\in[5;8)\) \\
4 & \(\dfrac58\) \\
5 & \(15\text{ м/с}\) \\
\bottomrule
\end{tabularx}
\end{table}

\vspace{10mm}

\noindent{\color{darkaccent}\textbf{Если ответ отличается, проверьте прежде всего:}}
\begin{enumerate}[leftmargin=8mm,itemsep=1.5mm]
  \item знак при переносе слагаемого;
  \item разворот знака при делении на отрицательное число;
  \item пересечение промежутков в системе;
  \item переворот второй дроби при делении;
  \item перевод единиц и порядок преобразования формулы.
\end{enumerate}

\vfill

\begin{center}
\small\textcolor{darkaccent}{Лёвин Артём Александрович}
\end{center}

\end{document}